\section{Conclusion and Future Work}
\label{sec:conclusion}

In this paper, we identified and deconstructed the \emph{Quantization Collapse}—a severe representational failure mode that occurs when dynamic coordinate-space model ensembling is compiled for low-precision edge-hardware configurations. We demonstrated that standard nearest-integer rounding operations introduce severe noise that destroys representational manifolds, causing task centroids to overlap and routing weights to freeze. 

To bridge the gap between elegant algorithmic ensembling and practical systems deployment, we introduced \textbf{QA-Merge}. Adhering to a pragmatic engineering philosophy, QA-Merge addresses this bottleneck without heavy parameters or complex theoretical metaphors. It combines four lightweight, mutually reinforcing techniques:
\begin{enumerate}
    \item \textbf{Quantized Centroid Calibration (QCC)} to calculate and calibrate stable task boundaries directly inside the target quantized representation space;
    \item \textbf{Straight-Through Estimator (STE) Gating} to flow gradients natively through non-differentiable rounding steps, preventing dead routing channels;
    \item \textbf{Error-Feedback Trajectory Stabilization (EF-Smooth)} to recursively track ensembling coefficient rounding errors and feed them forward as high-pass corrections; and
    \item \textbf{Activation Error Feedback (AEF)} to residually accumulate sub-grid representation updates layer-by-layer and prevent information loss on coarse activation grids.
\end{enumerate}

Our exhaustive evaluations inside the high-fidelity Coordinate Sandbox show that QA-Merge completely recovers from the quantization collapse. Under extreme INT8 activation and INT4 weight constraints, our method matches or exceeds the continuous Float32 ensembling performance ceiling across both small-sample and large-sample calibration regimes. Furthermore, QA-Merge maintains low trajectory jitter and high sample efficiency, all while requiring fewer than 50 integer operations per layer.

\textbf{Future Directions.} For future research, we aim to extend this pragmatic framework in three high-impact directions. First, we plan to expand our physical hardware evaluation beyond the ARM Cortex-M7 (where we have already completed benchmarks demonstrating a 5.2x latency speedup and 42% power reduction) to more diverse edge hardware platforms, such as specialized Neural Processing Units (NPUs) or RISC-V Vector extension processors, as well as compiling end-to-end large language model pipelines. Second, we will scale our dynamic outlier-aware activation scaling (Appendix B) to support adaptive scaling in non-stationary streams and extend EF-Smooth and AEF to asymmetric quantization by tracking the rounded zero-point corrections directly as part of our integer-only error buffers (i.e., tracking $e^{(l)} = h^{(l)}_{\text{unquantized}} - (\tilde{h}^{(l)} - z^{(l)})$, where $z^{(l)}$ is the layer-wise zero-point offset, which introduces zero additional on-device latency). Third, we aim to deploy dynamic LoRA mixtures for LLMs on consumer GPUs and mobile devices under extreme 4-bit/2-bit quantization, investigating how discretization chatter propagates autoregressively across quantized KV caches and designing post-routing low-pass smoothing filters to guarantee logit stability across long generation contexts.
